\documentclass[11pt]{article}

% ---------- Page & Font ----------
\usepackage[margin=2.5cm]{geometry}
\usepackage[utf8]{inputenc}
\usepackage[T1]{fontenc}
\usepackage{lmodern}
\usepackage{microtype}

% ---------- Math & Units ----------
\usepackage{amsmath, amssymb}
\usepackage{siunitx}
\sisetup{per-mode=symbol}

% ---------- Figures & Tables ----------
\usepackage{graphicx}
\usepackage{caption}
\usepackage{subcaption}
\usepackage{float}
\usepackage{booktabs}
\usepackage{array}
\usepackage[table]{xcolor}

% ---------- Code Listings (C, C++, Rust) ----------
\usepackage{listings}
\usepackage{xcolor}

\definecolor{codebg}{RGB}{245,245,245}
\definecolor{codegray}{RGB}{110,110,110}
\definecolor{codeblue}{RGB}{0,80,180}
\definecolor{codegreen}{RGB}{0,120,0}
\definecolor{codered}{RGB}{170,0,0}

\lstdefinelanguage{Rust}{
  keywords={fn, let, mut, pub, struct, enum, impl, trait, for, while, loop, if, else, match, return, use, mod, crate, self, Self, in, as, ref, move, unsafe, dyn, where, const, static, type},
  keywordstyle=\color{codeblue}\bfseries,
  ndkeywords={u8,u16,u32,u64,usize,i8,i16,i32,i64,isize,f32,f64,bool,char,str,String,Vec,Option,Result,Box},
  ndkeywordstyle=\color{codegreen},
  identifierstyle=\color{black},
  sensitive=true,
  comment=[l]{//},
  morecomment=[s]{/*}{*/},
  commentstyle=\color{codegray}\itshape,
  stringstyle=\color{codered},
  morestring=[b]",
}

\lstset{
  backgroundcolor=\color{codebg},
  basicstyle=\ttfamily\footnotesize,
  breaklines=true,
  breakatwhitespace=false,
  frame=single,
  rulecolor=\color{black!20},
  numbers=left,
  numberstyle=\tiny\color{codegray},
  numbersep=8pt,
  showstringspaces=false,
  tabsize=2,
  captionpos=b,
  keywordstyle=\color{codeblue}\bfseries,
  commentstyle=\color{codegray}\itshape,
  stringstyle=\color{codered},
}

% Convenience environments per language
\lstnewenvironment{ccode}[1][]{\lstset{language=C,#1}}{}
\lstnewenvironment{cppcode}[1][]{\lstset{language=C++,#1}}{}
\lstnewenvironment{rustcode}[1][]{\lstset{language=Rust,#1}}{}

% ---------- Misc ----------
\usepackage{hyperref}
\hypersetup{colorlinks=true, linkcolor=blue, citecolor=blue, urlcolor=blue}
\usepackage{enumitem}
\usepackage{xurl}

% ---------- Custom colored boxes for special sections ----------
\usepackage[most]{tcolorbox}

\newtcolorbox{hypothesisbox}{
  colback=blue!4, colframe=blue!50!black, title=Pre-Experiment Hypothesis,
  fonttitle=\bfseries, breakable
}
\newtcolorbox{failurebox}{
  colback=red!4, colframe=red!50!black, title=Negative / Failure Case,
  fonttitle=\bfseries, breakable
}
\newtcolorbox{anomalybox}{
  colback=yellow!8, colframe=orange!70!black, title=Honest Anomalies,
  fonttitle=\bfseries, breakable
}

% ---------- Title Info ----------
\title{\textbf{ECSE-6320 Advanced Computer Systems \\ Project 1 Report}}
\author{Cooper Werner \\\small Rensselaer Polytechnic Institute \\ \small wernec6@rpi.edu \\ \small RIN: 662058426}
\date{\today}

\begin{document}

\maketitle

\begin{abstract}
% One-paragraph summary: what was tested, what was expected, what was found.

\end{abstract}

\tableofcontents
\newpage

% =====================================================
\section{Introduction}
% Context, motivation, system under test, and objective of this experiment.

% =====================================================

The goal of this project is to analyze the performance of SIMD Instructions and how their performance is affected by the datasize,
size of the vector, memory alignment, data size, and compiler optimization.

To do this I will run multiple experiments across three different kernels:

\begin{enumerate}
	\item Vector Multiply and Accumulate: $\vec{y} = a \times \vec{x} + \vec{y} $
	\item Vector Elementwise Multiplication: $z_i  = x_i * y_i$
	\item 3-Element 1D Stencil: $y_i = a\times x_{i-1} + b \times x_{i} + c \times x_{i+1}$ 
\end{enumerate}

For each kernel I will adjust the size of the vector using different vectoriziation instructions (SSE, AVX, AVX512). Adjust memory alignment, element size (FP32 vs FP64),
number of elements provided, and compiler flags for autovectorization. 

Performance will be evaluated using the Linux \texttt{timespec} to measure execution time of the vector execution while isolating overhead from I/O operations,
and verification.

\section{Pre-Experiment Hypotheses}

\begin{hypothesisbox}
\textbf{Hypothesis:} \\
% State the falsifiable prediction before running anything.

The compiler's autovectorization will produce code that runs significantly faster than the serial version. Since the compiler doesn't have the same
knowledge as my code as I do, I think my SIMD implementation will be at least the same speed as the autovectorized version or a bit faster.

\textbf{Reasoning:} \\
% Why do you expect this outcome? Reference architecture, prior results, theory.
The compiler will target the largest vectors that are available on my CPU since it will be provided with the \texttt{-march=native} compiler flag.
Although my kernels are fairly simple vector operations, since I have better knowledge of my data layout, I think that my handwritten vectoriziation
will perform better than the compiler's vectoriziation.

\textbf{Predicted quantitative outcome:} \\
\begin{center}
\captionof{table}{Predicted SIMD speedup relative to scalar baseline}
\label{tab:hypothesis-simd}
\begin{tabular}{lccc}
\toprule
\textbf{Vectoriziation Strategy} & \textbf{FP32 Speedup from Serial} & \textbf{FP64 Speedup from Serial} \\
\midrule
Serial & $1x$ & $1x$ \\
Compiler Autovectorization & $10.5x$ & $5.2x$ \\
Handwritten Vectoriziation & $11x$ & $5.6x$ \\
\bottomrule
\end{tabular}
\end{center}
\end{hypothesisbox}

\begin{hypothesisbox}
\textbf{Hypothesis:} \\
% State the falsifiable prediction before running anything.

Shifting the memory alignment off of the 64 Byte alignment that is preferred by AVX512 will significantly affect the performance of the SIMD instructions. 
The instructions will likely still run, however they will not run efficiently and may even be slower than the scalar counterpart.

\textbf{Reasoning:} \\
% Why do you expect this outcome? Reference architecture, prior results, theory.

CPUs are most efficient when the data is alligned to the memory boundary since CPUs will fetch memory multiple bytes at a time. If the memory is shifted off 
the memory alignment boundary it will have to do multiple data fetches to fill the buffers. 
This will increase the number of pipeline stalls and reduce instruction throuhput.

\textbf{Predicted quantitative outcome:} \\
% e.g., "Expect L2 miss rate to decrease by \SI{15}{\percent} \pm \SI{3}{\percent}"
\begin{center}
\captionof{table}{Predicted SIMD Slowdown relative to SIMD with 0 byte offset}
\label{tab:hypothesis-simd}
\begin{tabular}{lccc}
\toprule
\textbf{Offset Bytes} & \textbf{FP32 Slowdown} & \textbf{FP64 Slowdown} \\
\midrule
 1  & $\sim$0.5$\times$ & $\sim$0.25$\times$ \\
 2 & $\sim$0.5$\times$ & $\sim$0.25$\times$ \\
 3  & $\sim$0.5$\times$  & $\sim$0.25$\times$ \\
 4  & $\sim$1$\times$  & $\sim$0.25$\times$ \\
 5  & $\sim$0.5$\times$  & $\sim$0.25$\times$ \\
 6  & $\sim$0.5$\times$  & $\sim$0.25$\times$ \\
 7  & $\sim$0.5$\times$  & $\sim$0.25$\times$ \\
\bottomrule
\end{tabular}
\end{center}
\end{hypothesisbox}

% =====================================================
\section{Correctness Validation}
% How do you know the system under test / measurement apparatus is functioning correctly
% BEFORE trusting the results? E.g. reference outputs, checksums, sanity benchmarks,
% comparison against known-good baseline, unit tests.

\subsection{Validation Method}


\subsection{Validation Results}

\begin{table}[H]
\centering
\caption{Correctness validation summary}
\label{tab:validation}
\begin{tabular}{@{}lll@{}}
\toprule
Check & Expected & Observed \\
\midrule
 &  &  \\
 &  &  \\
\bottomrule
\end{tabular}
\end{table}

% =====================================================
\section{Measurement Discipline}
% Describe methodology: isolation (core pinning, turbo/frequency scaling disabled, etc.),
% number of trials, warm-up runs, statistical treatment (mean/median/stdev/CI),
% timer resolution, environmental controls.

\subsection{Experimental Setup}
\begin{itemize}[leftmargin=1.5em]
    \item \textbf{Hardware:} CPU: x86\_64 AMD Ryzen AI 7 350 w/ Radeon 860M (SIMD Support: SSE, SSE2, AVX, AVX2, AVX-512), Memory: 32GB LPDDR5 @ 5600MT/s
    \item \textbf{OS / Kernel:} Pop\_OS! 24.04 Running Linux 7.0.11-76070011-generic
    \item \textbf{Compiler \& Flags:} GCC 16.2.0, Flags: -O2 -march=znver4 -ftree-vectorize (disabled for serial experiments)
    \item \textbf{Isolation techniques:} Ran on CPU0 with 100\% affinity, and -20 niceness (ensure no context switching)
    \item \textbf{Trials per condition:} 5 trials per condition
    \item \textbf{Timer / Instrumentation:} clock\_gettime(CLOCK\_MONOTONIC)
\end{itemize}

\subsection{Statistical Treatment}
% Mean, stdev, confidence intervals, outlier rejection criteria (and whether/why applied).

% =====================================================
\section{Results}

\subsection{Summary Figure}
\begin{figure}[H]
    \centering
    % \includegraphics[width=0.85\textwidth]{figures/placeholder.png}
    \caption{Caption describing the primary result.}
    \label{fig:main-result}
\end{figure}

\subsection{Secondary Figures}
\begin{figure}[H]
    \centering
    \begin{subfigure}[b]{0.48\textwidth}
        \centering
        % \includegraphics[width=\textwidth]{figures/placeholder1.png}
        \caption{}
        \label{fig:sub1}
    \end{subfigure}
    \hfill
    \begin{subfigure}[b]{0.48\textwidth}
        \centering
        % \includegraphics[width=\textwidth]{figures/placeholder2.png}
        \caption{}
        \label{fig:sub2}
    \end{subfigure}
    \caption{Side-by-side comparison of \dots}
    \label{fig:sub-both}
\end{figure}

\subsection{Results Table}
\begin{table}[H]
\centering
\caption{Quantitative results}
\label{tab:results}
\begin{tabular}{@{}lccc@{}}
\toprule
Configuration & Mean & Std.\ Dev. & N \\
\midrule
 &  &  &  \\
 &  &  &  \\
\bottomrule
\end{tabular}
\end{table}

\subsection{Relevant Code}

% --- C example ---
\begin{ccode}[caption={Benchmark kernel (C)}, label={lst:c-example}]
int main(void) {
    // ...
    return 0;
}
\end{ccode}

% --- C++ example ---
\begin{cppcode}[caption={Measurement harness (C++)}, label={lst:cpp-example}]
#include <chrono>

int main() {
    auto start = std::chrono::high_resolution_clock::now();
    // ...
}
\end{cppcode}

% --- Rust example ---
\begin{rustcode}[caption={Benchmark harness (Rust)}, label={lst:rust-example}]
fn main() {
    let start = std::time::Instant::now();
    // ...
}
\end{rustcode}

% =====================================================
\section{Negative / Failure Case}

\begin{failurebox}
\textbf{What failed:} \\
% Configuration, workload, or condition that did NOT behave as expected or hypothesized.

\textbf{Evidence:} \\
% Data, logs, or figure reference showing the failure.

\textbf{Root cause (if known):} \\

\end{failurebox}

% =====================================================
\section{Honest Anomalies}

\begin{anomalybox}
% Report unexplained results honestly, even if they don't fit the narrative.
% Include: what was observed, what was expected instead, and what was (or wasn't) ruled out.

\textbf{Anomaly 1:} \\


\textbf{Anomaly 2:} \\

\end{anomalybox}

% =====================================================
\section{Discussion}
% Tie hypothesis, validation, results, and anomalies together. Was the hypothesis supported,
% refuted, or inconclusive? What would you do differently next time?

% =====================================================
\section{Conclusion}


% =====================================================
\appendix
\section{Additional Data / Raw Logs}
% Optional: full result dumps, extra plots, raw perf output, etc.

\end{document}